\documentclass[11pt,a4paper]{article}
\usepackage[T1]{fontenc}
\usepackage[utf8]{inputenc}
\usepackage{lmodern,microtype}
\usepackage[margin=24mm,headheight=15pt]{geometry}
\usepackage{amsmath,amssymb,mathtools}
\usepackage{booktabs,longtable,tabularx,array}
\usepackage{xcolor,listings,enumitem}
\usepackage{xurl,hyperref,fancyhdr}
\definecolor{accent}{HTML}{24465B}
\definecolor{lightgray}{gray}{0.96}
\hypersetup{colorlinks=true,linkcolor=accent,urlcolor=accent,citecolor=accent,
 pdftitle={Asymptotic: Correctness, Contracts, and Engineering Audit},
 pdfauthor={Repository audit prepared for Vladimir Reshetnikov}}
\urlstyle{same}
\pagestyle{fancy}
\fancyhf{}
\fancyhead[L]{\small\sffamily Asymptotic\enspace|\enspace Repository audit}
\fancyhead[R]{\small\sffamily 9 September 2026}
\fancyfoot[C]{\small\thepage}
\setlength{\parskip}{5pt}
\setlength{\parindent}{0pt}
\setlength{\emergencystretch}{3em}
\renewcommand{\arraystretch}{1.18}
\newcolumntype{P}[1]{>{\raggedright\arraybackslash}p{#1}}
\setcounter{tocdepth}{2}
\newcommand{\code}[1]{\texttt{\detokenize{#1}}}
\newcommand{\R}{\mathbb R}
\newcommand{\N}{\mathbb N}
\newcommand{\OO}{\mathcal O}
\newcommand{\SHA}{\code{07a9781212beb2eeb9ff16aa625b50ac27974078}}
\newcommand{\finding}[3]{\subsection{#1: #2}\textbf{Classification.} #3\par}
\lstset{basicstyle=\ttfamily\footnotesize,breaklines=true,breakatwhitespace=false,
 columns=fullflexible,keepspaces=true,showstringspaces=false,
 backgroundcolor=\color{lightgray},frame=single,rulecolor=\color{black!15},
 framerule=0.3pt,xleftmargin=3pt,xrightmargin=3pt,
 aboveskip=9pt,belowskip=9pt,tabsize=2}
\begin{document}
\begin{titlepage}
\thispagestyle{empty}
\vspace*{16mm}
{\sffamily\large\color{accent} TECHNICAL REPOSITORY REVIEW}\par
\vspace{9mm}
{\sffamily\bfseries\Huge Asymptotic}\par
\vspace{4mm}
{\sffamily\LARGE Correctness, Contracts,\par and Engineering Audit}\par
\vspace{9mm}
{\large Comparison with built-in Wolfram Language capabilities,\par source-derived findings, and a development plan}\par
\vspace{16mm}
\begin{tabular}{@{}ll@{}}
Repository & \code{VladimirReshetnikov/Asymptotic}\\
Reviewed version & 1.8.0\\
Pinned revision & \SHA\\
Review date & 9 September 2026\\
Deliverable & Article, proposed patches, tests, and evidence\\
\end{tabular}
\vfill
\begin{minipage}{0.94\linewidth}
\textbf{Evidence boundary.} This is a source audit with independently executed
mathematical and Python-tool checks. The connected native Wolfram service
returned errors, and no local native kernel was available. Consequently,
this review does \emph{not} claim to have executed the package, reproduced
its reported Wolfram tests, measured its runtime, or validated the proposed
Wolfram changes in a native kernel. Native test and benchmark harnesses are
included for that purpose.
\end{minipage}
\vspace{10mm}\par
{\sffamily Prepared for Vladimir Reshetnikov}
\end{titlepage}

\begin{abstract}
The repository is a substantial real-asymptotic computation library, not
merely a wrapper around \code{InverseSeries}. Its strongest contribution is
an explicit computational representation of scales, branches, complete
coefficient blocks, transported remainders, and reusable inverse
computations. It complements Wolfram Language's considerably broader
symbolic-analysis system. This review identifies two concrete hazards in
the optional \code{SeriesData} export path, an observable-power option
precedence problem, and a source-derived example of under-refinement caused
by a fixed precision margin. It also examines expensive canonicalization,
resource budgeting, native-series ingestion, validation provenance,
certificate scope, API consistency, and migration costs. The recommendations
prioritize preserving mathematical contracts and making achieved precision
observable before adding more function families. The accompanying artifacts
separate executed independent checks from unexecuted native regression tests
and conservative patch proposals.
\end{abstract}
\tableofcontents
\newpage

\section{Executive assessment}
\label{sec:executive}
\textbf{The project is worth developing as a contract-aware layer above the
Wolfram kernel.} Its value is not that Wolfram Language has no asymptotics,
no fractional series, or no inverse expansion machinery. Those claims would
be inaccurate. Its value is that it packages a particular mathematical
workflow: select a real germ, work in an explicit scale, compute complete
blocks, preserve an error model, and retain enough information to inspect or
refine the result. The source implements ordinary power--log jets,
logarithmic and exact-core reductions, Fourier-polynomial coefficients,
finite flat sectors, special-function adapters, arithmetic, refinement, and
several distinct kinds of validation.\cite{readme,core,seriesops,fourier,sourcecoords}

The library also contains safeguards that should not be sacrificed for
apparent convenience. It does not equate a small numerical residual with a
proof. The interval certificate explicitly concerns a stored, explicit
function on a checked interval, rather than every function compatible with
an unspecified asymptotic remainder. Differentiation requires a derivative
contract. Fourier inverses reject an oscillatory leading block rather than
inventing an eventual sign. Special-function input models distinguish finite
Poincare information from convergent exact data.\cite{certificates,seriesops,fourier,special}

The highest-priority findings concern \emph{interfaces between otherwise
reasonable components}:
\begin{enumerate}[leftmargin=*,itemsep=3pt]
\item A sparse expansion is eagerly converted to a dense rational-lattice
\code{SeriesData}. Two retained blocks can request a billion coefficient
slots. Neither the sparse term count nor the user's \code{"MaxTerms"}
budget bounds this allocation.
\item Both classical-series exporters omit the logarithmic degree of the
remainder. The rich object can correctly describe
$\OO(x|\log x|)$ while the exported object no longer carries that contract.
\item The inverse-object coefficient wrapper prepends an inherited
\code{"Power"} option before explicit caller options. The explicit setting
is therefore silently shadowed.
\item Derived-operation refinement uses a fixed guard margin rather than
propagating precision demands through valuations. A large pole supplies a
simple example in which a valid returned approximation can remain far below
the requested precision.
\end{enumerate}
The first three findings are visible directly in the relevant source paths;
the fourth includes an explicit mathematical trace, but all native
reproductions remain pending. Detailed classifications and examples appear
in Sections~\ref{sec:registry}--\ref{sec:details}.\cite{core,seriesops,optionvalue}

The development order should be: contain the export hazards; make native
validation a release gate; introduce explicit achieved-precision and
capability descriptions; improve operation-level precision planning; then
expand analytic scope. Replacing the mature kernel's general symbolic
solvers would be a distraction. Adapting their output into carefully checked
contracts is a more credible strategy.

\section{Snapshot, method, and limits of this audit}
\label{sec:method}
\subsection{The object actually reviewed}
All repository conclusions refer to revision \SHA. The paclet manifest calls
the package \code{AsymptoticInverse}, gives version 1.8.0, and requires
\code{WolframVersion -> "15.0+"}. The tree contains 38 kernel \code{.wl}
modules, in addition to \code{init.m}. The repository-root standalone
\code{AsymptoticInverse.wl} is a generated distribution; the modular kernel
is canonical. Historical reports and their separate implementations are not
alternative current production entry points.\cite{tree,paclet,development}

This distinction matters. A defect in an archived proposal is not a defect
in the maintained package. Conversely, a successful historical manifest does
not establish correctness of every current module. The report therefore
uses pinned source links and function names, not moving \code{main} links,
for its implementation findings. The archive contains a module inventory,
a source-coverage manifest, and an explicit native-execution status file.

\subsection{Evidence levels}
\begin{longtable}{@{}P{0.10\linewidth}P{0.25\linewidth}P{0.58\linewidth}@{}}
\toprule
Label & Meaning & What it does and does not establish\\\midrule
\endhead
S & Source observation & A control path, option list, allocation, or data
field is present in the reviewed text. This is not a measured runtime.\\
D & Mathematical derivation & A stated consequence follows from an explicit
model or contract. Its assumptions are supplied.\\
E & Executed independent check & The supplied Python/SymPy or Python-tool
check was run in this environment. It does not execute the Wolfram package.\\
R & Repository-reported result & An upstream validation record is being
reported, not independently reproduced.\\
U & Native reproduction pending & A Wolfram test or benchmark is supplied,
but no native result is claimed.\\
P & Proposal or risk & A recommended design, engineering risk, or audit
candidate, not an established wrong public answer.\\
\bottomrule
\end{longtable}
A single finding can have several labels. For example, the dense allocation
has a source-observed allocation, an exact derivation of its length, an
independently executed size calculation, and an unexecuted native fixture.
This is stronger than a vague performance suspicion, but it still is not a
measured memory benchmark.

\subsection{Inspection and execution scope}
The audit followed the public entry points through the core jet algebra,
forward and inverse construction, result formatting, classical-series
export, residuals, coefficient queries, series arithmetic and refinement,
source-coordinate reconstruction, Fourier inverses, native special-function
ingestion, Gamma inverse construction, and elementary interval certificates.
Packaging metadata, the CI workflow, the full-suite runner, development
notes, and current validation summaries were also inspected. Some other
family modules were inventoried and assessed through their dispatch paths
and repository descriptions rather than exhaustively line-audited.
Appendix~\ref{sec:coverage} records that distinction.

The native Wolfram connector's context and evaluation calls returned service
failures. No native Wolfram executable was present locally. The independent
mathematical checks ran with Python 3.13.5 and SymPy 1.14.0. There are 16
successful mathematical checks and seven successful Python-tool tests in the
included evidence. These counts must not be added to the repository's native
pass counts as though they were one suite.

\subsection{Current upstream validation is substantial but scoped}
The latest 1.8.0 section of the validation record reports 305 passing tests
from an explicitly selected group of 18 files, zero failures, builder and
documentation checks, and isolated standalone acceptance checks. It also
explicitly says that the full package suite was not run at that milestone.
Its operation timings compare specific earlier and later package fixtures;
they are not comparisons with Wolfram built-ins.\cite{validation}

The older \code{review-manifest.json} records a different baseline commit,
four kernel files, and 130 successes. It is a historical artifact, not a
current all-module attestation. The development documentation already warns
readers to preserve that historical scope. The recommendation here is to
make the current release attestation equally easy to identify and verify,
not to erase the historical record.\cite{reviewmanifest,development}

\section{What the package implements}
\label{sec:architecture}
\subsection{A scale-aware result rather than only a formula}
The central public head is \code{GeneralizedSeries}. An object retains a
finite expression together with properties such as its variable, approach,
terms, remainder, and provenance; specialized kinds add different data.
\code{Normal[s]} deliberately returns only the finite expression. Standard
and traditional formatting hide the head but retain the original object
inside an interpretation. The formatting code uses held forms and read-only
interpretation boxes to avoid executing payloads or retaining stale metadata
after a visual coefficient edit. These are useful notebook-safety choices.
\cite{core}

The ordinary computational representation is a sparse jet
\[
 J=(T,\rho,d),\qquad
 T=\{(\alpha_j,P_j(L))\},\qquad L=\log u,
\]
representing a retained sum of $u^{\alpha_j}P_j(\log u)$ and an error class
\[
 \OO\!\left(u^\rho(1+|\log u|)^d\right),\qquad u\downarrow0.
\]
An infinite precision marker denotes absence of an unknown tail in the
internal representation. Equal weights must be merged before a block is
counted or a truncation frontier is selected. The code canonically handles
rational and algebraic weights, attempts exact treatment of further real
constants, and refuses an undecidable order rather than intentionally
sorting approximate exponents.\cite{core}

\subsection{Endpoint and branch normalization}
For the ordinary real approaches, the local coordinate is one of
\[
 u=x-x_0,\quad u=x_0-x,\quad u=1/x,\quad u=-1/x.
\]
Each is positive on the selected side and tends to zero. This allows one
local algebra to serve finite endpoints and both infinities. Additional
source-coordinate modules recognize certain exponential and logarithmic
charts, reconstruct the requested source observable, and reject an input
remainder that has not first been transported into the chart. They also
retain exact chart inverses when exponentiating an inexact intermediate
would amplify its error unacceptably.\cite{core,sourcecoords}

A subtle but important API convention is that a nonunit \code{"Power"} at a
finite source endpoint describes a local source observable. The formula is
not automatically the global power $x^r$ after an arbitrary translation.
The source code's coordinate sign and endpoint reconstruction must therefore
remain part of any coefficient or refinement API. A generic wrapper that
forgets this distinction can be wrong even when its formal series algebra
is correct.\cite{core,refinement}

\subsection{Implemented scale families and their boundaries}
\begin{longtable}{@{}P{0.24\linewidth}P{0.42\linewidth}P{0.27\linewidth}@{}}
\toprule
Area & Existing capability & Important boundary\\\midrule
\endhead
Ordinary power--log & Exact sparse weights, logarithmic-polynomial blocks,
Lagrange/Newton-related inversion, arithmetic and refinement & Not every
symbolic exponent order is decidable.\\
Lambert and logarithmic reductions & Leading logarithmic/exponential cores
and logarithmic coordinate scales & Cutoff meanings differ from an ordinary
target-power cutoff.\\
Fourier coefficients & Finite modes in the logarithmic coordinate, complete
frequency grouping and bounded coefficient envelopes & An oscillatory
leading core is not given an unproved eventual sign.\\
Flat and exponential structure & Finite flat sectors and operations,
retained exponential carriers, and composite error envelopes & Not a claim
of a universal transseries or resurgent-analysis engine.\\
Gamma and Barnes families & Dedicated positive-argument asymptotic models,
Lambert-based inverse cores, residual and numerical checks & Finite
Poincare information is not convergent exact input.\\
Other special functions & Explicit adapters plus native structured-series
ingestion, parameter and real-domain checks & Recognizing a function head
does not prove support for every parameter and endpoint.\\
Certificates & Elementary exact rational interval root enclosures for
admitted inverse kinds & Not all specialized inverse kinds, arbitrary
special-function expressions, or unspecified remainder families.\\
\bottomrule
\end{longtable}
This inventory is grounded in the current public description and the cited
implementation paths. It would be misleading to propose Fourier support,
Gamma inverses, finite flat-sector arithmetic, or retained refinement state
as if they were absent. The opportunities lie in making their contracts
more uniform and extending their interoperability.\cite{readme,fourier,native,gamma,special,certificates,refinement}

\subsection{Reusable computation is already present}
The refinement implementation retains coefficient prefixes and boundaries
or Newton unit jets, verifies compatibility, and includes a signature of the
model, observable, assumptions, function, variables, endpoint, and direction.
It uses value-like retained states rather than a global cache. Replayed
operations are explicitly distinguished from reused computations in
\code{RefinementStatistics}. These are good foundations for both performance
and reproducibility. The remaining issue is not simply ``add caching''; it
is making precision planning, cache accounting, and replay coverage as
systematic as the stored-state machinery.\cite{refinement,seriesops}

\section{Comparison with built-in Wolfram Language}
\label{sec:builtins}
\subsection{The correct comparison is complementary, not categorical}
Wolfram's documented \code{Asymptotic} supports more than elementary power
series: it infers scales, accepts finite or infinite real or complex
expansion points, and can operate on expressions involving integrals,
transforms, and differential-equation solutions. Its options include
assumptions, conditions, precision and performance controls, and
\code{SeriesTermGoal}. Accordingly, inability of one \code{SeriesData}
representation to encode a particular sparse scale does not establish that
Wolfram Language cannot compute any approximation in that scale.
\cite{asymptotic}

\code{AsymptoticSolve} is a particularly relevant baseline. It supports
implicit equations and systems, can select solutions passing through a
specified point, has a real-valued-solution form, and documents real and
complex approach directions. Its approximation order is not generally the
polynomial degree. The package's inverse routines address a narrower but
more inspectable workflow, rather than replacing all of these facilities.
\cite{asymptoticsolve}

\subsection{Capability matrix}
\begin{longtable}{@{}P{0.23\linewidth}P{0.34\linewidth}P{0.36\linewidth}@{}}
\toprule
Task & Native baseline & Package contribution or limitation\\\midrule
\endhead
Taylor, Laurent, and Puiseux & \code{Series}, \code{SeriesData} & Sparse explicit
blocks and richer branch/remainder provenance; native tools are the first
baseline on conventional cases.\\
Ordinary reversion & \code{InverseSeries} & Complete power--log and exact
weight model machinery, observable powers, and reusable computation.\\
Implicit asymptotic solutions & \code{AsymptoticSolve} & Specialized
single-real-branch inverse construction with retained model and checks;
less general than native systems solving.\\
Composition & \code{ComposeSeries} and series arithmetic & Explicit
transport of heterogeneous precision and carrier information in admitted
scales.\\
Coefficient extraction & \code{SeriesCoefficient} & Model multi-index
coefficients and Fourier-specific queries; these are not interchangeable
with coefficients indexed only by a final power.\\
General asymptotics & \code{Asymptotic} & Durable typed results, inspection,
refinement, and explicit contracts for the supported families.\\
Logarithmic inversion & \code{ProductLog}, symbolic solving and asymptotics
& Explicit branch and scale bookkeeping around recognized cores.\\
Integrals & \code{AsymptoticIntegrate} & The reviewed package is not a
general replacement for asymptotic integration.\\
Discrete problems & \code{DiscreteAsymptotic} and related native facilities
& Dedicated inverse/continuous-germ contracts are a different focus.\\
Differential equations & \code{AsymptoticDSolveValue} & No comparable
general differential-equation solver is established by this audit.\\
Asymptotic comparisons & \code{AsymptoticLess} and related relations &
Explicit exponent, logarithmic-degree, envelope, and coordinate metadata
can support mechanical contract comparisons.\\
Certified numerical inversion & Native numerical solvers and symbolic
analysis are useful ingredients & An elementary rational-interval
certificate with a stated local scope is an additional package facility;
it is distinct from merely calling \code{FindRoot}.\\
\bottomrule
\end{longtable}
Native rows refer to the corresponding official documentation; package
claims refer to the inspected implementation and current overview.
\cite{series,seriesdata,inverseseries,composeseries,seriescoefficient,asymptotic,
asymptoticsolve,productlog,asymptoticintegrate,discrete,asymptoticdsolve,
asymptoticless,readme,certificates}

\subsection{Representation versus mathematical ability}
\code{SeriesData} represents coefficients on a regular exponent lattice,
with exponents determined by integer positions and a denominator. It is
well suited to conventional rational-power series. A sparse collection such
as $1+x^{\sqrt2}+x^{\sqrt3}$ has a different natural representation. Even for
rational powers, a denominator can make a dense lattice inefficient.
Neither observation is a blanket statement about the reach of the native
\code{Series} or \code{Asymptotic} algorithms.\cite{seriesdata,series}

\code{InverseSeries} and \code{ComposeSeries} provide native series
reversion and composition, including conventional fractional-power
situations. They should be used as differential baselines wherever the
input and output contracts genuinely overlap. \code{SeriesCoefficient}
also offers direct coefficient queries from series or expressions. The
package's multi-index query has a different meaning: several model
multi-indices can contribute to the same final exponent and must be added
before one obtains its complete coefficient block.
\cite{inverseseries,composeseries,seriescoefficient,core}

\subsection{Order and term counts must be normalized before comparison}
For the ordinary package path, a cutoff $H$ is exclusive: retain blocks with
weight strictly below $H$. A native call such as \code{Series[f,{x,0,n}]}
uses the native series order convention. An \code{AsymptoticSolve} order
can count approximation stages rather than polynomial powers. Equal
integers in these calls therefore do not establish equal work or equal
accuracy.\cite{core,series,asymptoticsolve}

A meaningful comparison records the local coordinate, selected branch,
retained expression, first omitted scale, logarithmic degree, and whether
validation work was included. A two-line native expression should not be
called ``faster'' than a package object that additionally constructs a
certificate without separately accounting for that certificate. Conversely,
a large metadata object is not inherently evidence of a more accurate
approximation.

\subsection{Recommended user workflow}
Use native \code{Series} and \code{InverseSeries} for standard local series
when they satisfy the task. Use \code{Asymptotic} or
\code{AsymptoticSolve} for broad symbolic exploration and implicit problems.
Use this package when complete nonstandard blocks, a recorded real branch,
precision-preserving operations, or repeatable inverse refinement are the
main requirement. Then use exact identities or a separately conditioned
numerical check to validate the overlap. These are workflow recommendations,
not measured claims of universal superiority.

\section{Mathematical contracts that the software must preserve}
\label{sec:contracts}
\subsection{Ordinary inverse model and complete coefficients}
Consider, on a specified real germ,
\begin{equation}
 v=a u^p\left(1+\sum_{j=1}^{m}u^{\delta_j}B_j(\log u)\right),
 \qquad a\ne0,\quad p\ne0,\quad \delta_j>0,
 \label{eq:model}
\end{equation}
where $u>0$ tends to zero and the branch satisfies $v/a>0$. Set
$z=(v/a)^{1/p}$ and $L=\log z$. For a fixed local observable $u^r$, the
formal expansion has the form
\begin{equation}
 u(z)^r=z^r\sum_{\boldsymbol{k}\in\N^m}
 z^{\boldsymbol{k}\cdot\boldsymbol\delta}C_{\boldsymbol{k}}(L).
 \label{eq:inverse}
\end{equation}
The core's Euler-operator formula, also reflected in the Fourier
implementation, can be written as follows. If $n=|\boldsymbol{k}|>0$,
$\Delta=\boldsymbol{k}\cdot\boldsymbol\delta$, and $D=d/dL$, then
\begin{equation}
 C_{\boldsymbol{k}}(L)=
 \frac{(-1)^n r}{p^n\prod_j k_j!}
 \prod_{q=1}^{n-1}\left(D+r+\Delta+qp\right)
 \prod_j B_j(L)^{k_j},
 \qquad C_{\boldsymbol0}=1.
 \label{eq:coefficient}
\end{equation}
The operators commute because they are polynomials in $D$. This formula
explains both the exact coefficient feature and why logarithmic degree
must be tracked separately from power.\cite{core,fourier}

For completeness, introduce independent perturbation markers in
\eqref{eq:model}. The observable version of Lagrange--Buermann applies
$(n-1)$ target derivatives to the product of the core observable derivative
and $n$ perturbation factors. A target derivative becomes a factor
$p^{-1}z^{-p}(D+\lambda)$ when applied to $z^\lambda P(\log z)$.
Collecting powers and selecting a multi-index yields the multinomial factor
$n!/\prod k_j!$, which cancels the $n!$ in the formula. The remaining
factors are exactly those in \eqref{eq:coefficient}. This derivation is
formal coefficient algebra; applicability to an unknown real function still
requires the branch and remainder hypotheses discussed below.

For $f(u)=u+u^2\log u$, independently solving the composition residual gives
\begin{align}
 f^{-1}(y)={}&y-y^2L+y^3(2L^2+L)\notag\\
 &-y^4\left(5L^3+\frac{11}{2}L^2+L\right)\notag\\
 &+y^5\left(14L^4+\frac{73}{3}L^3+\frac{21}{2}L^2+L\right)
 +\cdots,\qquad L=\log y.
 \label{eq:logoracle}
\end{align}
The included Python check obtains these coefficients by introducing one
unknown coefficient at a time and cancelling the residual, rather than
copying the package's Euler implementation. It also verifies zero residual
through the displayed degree. This is an independent algebraic oracle,
not an execution result from the package.

\subsection{Truncation is about complete layers}
Suppose different multi-indices satisfy
$\boldsymbol{k}\cdot\boldsymbol\delta=
\boldsymbol{l}\cdot\boldsymbol\delta$. Their contributions belong to the
same block. Truncating before all such contributions are added can leave a
wrong coefficient or an incorrect nonzero frontier after cancellation.
Even simple algebraic descriptions can conceal equality:
$\sqrt8=2\sqrt2$. The ordinary and Fourier implementations already devote
machinery to grouping these layers. The recommended optimizations must
preserve that invariant rather than reverting to a first-$n$-indices
heuristic.\cite{core,fourier,refinement}

\subsection{Remainders have more structure than an exponent}
At a fixed positive $u$, the notation
\[
 R(u)=\OO\!\left(u^\rho(1+|\log u|)^d\right)
\]
asserts an eventual bound with some constant and some neighborhood. It is not
itself a computable error bound at that $u$. It also cannot generally be
replaced by $\OO(u^\rho)$ when $d>0$. For every fixed $\varepsilon>0$,
\begin{equation}
 u^\varepsilon(1+|\log u|)^d\longrightarrow0,
 \qquad
 u^\rho(1+|\log u|)^d=\OO(u^{\rho-\varepsilon}),
 \label{eq:epsilon}
\end{equation}
but the loss of $\varepsilon$ is real information loss. This observation is
the mathematical basis for the conservative native-ingress policy examined
later, and for finding A02.

\subsection{Transporting source error through an inverse}
Assume the source model has an error $\OO(u^\rho L_u^d)$, where
$L_u=1+|\log u|$, and that the actual derivative satisfies the matching
nondegenerate branch behavior $f'(u)\sim apu^{p-1}$. The inverse-coordinate
error is then, by local mean-value conditioning,
\[
 \OO(u^{\rho-p+1}L_u^d).
\]
For the observable $u^r$, the corresponding error is
\begin{equation}
 \OO(u^{\rho-p+r}L_u^d).
 \label{eq:inverseerror}
\end{equation}
In the positive target error coordinate this gives exponent
$(\rho-p+r)/|p|$. At a source infinity, the internal observable exponent is
changed as required by $x=\pm1/u$. This explains why increasing the number
of formally inverted coefficients cannot exceed the information supplied
by the source remainder.\cite{core,fourier,refinement}

The derivative hypothesis is not cosmetic. An arbitrary function can have
a small magnitude remainder and a very large derivative remainder.
Consequently, declarations of \code{"InputRemainder"} must retain their
hypothesis status, including any matching derivative assumption, rather
than being displayed as a theorem automatically proved about the original
function.

\subsection{Differentiation, exponentiation, and conditioning}
Take $R(x)=x^2\sin(x^{-3})$. Then $R=\OO(x^2)$, but
\[
 R'(x)=2x\sin(x^{-3})-3x^{-2}\cos(x^{-3}).
\]
Along $x_n=(2\pi n)^{-1/3}$, $R'(x_n)/x_n=-6\pi n$, so
$R'\notin\OO(x)$. The package's refusal to differentiate a bare magnitude
contract is therefore correct.\cite{seriesops}

Similarly, if $h=H+R$, then $e^h=e^H e^R$. A small \emph{relative} error in
$h$ does not imply a small relative error in $e^h$ when $h$ diverges. One
needs an appropriate absolute control on $R$, or an exact retained inner
expression. The source-chart reconstruction and exponentiation paths
already encode this distinction; future generalized composition should use
it as an invariant.\cite{seriesops,sourcecoords}

\section{Findings registry and priorities}
\label{sec:registry}
Severity is relative to this library's purpose. A high-priority finding can
be a resource-exhaustion or contract-preservation defect without being a
wrong coefficient. ``Pending native'' means exactly that; source inspection
is not presented as an observed Wolfram session.

\begin{longtable}{@{}P{0.065\linewidth}P{0.35\linewidth}P{0.12\linewidth}P{0.36\linewidth}@{}}
\toprule
ID & Finding & Priority & Evidence and disposition\\\midrule
\endhead
A01 & Eager dense classical-series export can overwhelm a sparse computation
& High & S, D, E, U. Contain immediately; make conversion lazy.\\
A02 & Classical-series export loses logarithmic remainder metadata
& High & S, D, E, U. Reject or explicitly describe lossy conversion.\\
A03 & Inherited observable power silently shadows an explicit option
& Medium & S, D, U. Decide and enforce the public precedence contract.\\
A04 & A fixed replay margin can fail to achieve requested refinement
& Medium & S, D, E, U. Propagate precision demands through valuations.\\
A05 & Coefficient normalization can invoke uncontrolled integer factorization
& Medium & S, P, U. Budget and defer nonessential normalization.\\
A06 & One resource option measures several incompatible kinds of work
& Medium & S, P. Separate budgets and report the exhausted dimension.\\
A07 & Checked-in CI does not gate native mathematical correctness
& High & S, R, P. Add a pinned, native release attestation.\\
A08 & The legacy full-suite runner can accept an empty suite
& Medium & S, P, U. Reuse the stronger focused-runner infrastructure.\\
A09 & Cutoff and term-goal meanings change with the selected scale
& Medium & S, P. Expose a typed request and an execution plan.\\
A10 & Convenient numeric evaluation does not check domain or error
& Medium & S, P. Add a distinct checked-evaluation interface.\\
A11 & Result schemas, capabilities, and saved-head migration are implicit
& Medium & S, R, P. Version data and provide capability queries.\\
A12 & Native conventions are only partially accepted at the boundary
& Low--medium & S, P. Normalize direction aliases and document order differences.\\
A13 & Certificate coverage and the \code{Automatic} interval default need clearer discovery
& Low--medium & S, P. Improve required-argument diagnostics and capability descriptions.\\
A14 & Paclet documentation and license metadata need distribution polish
& Low & S, P. Add native help integration and align license identifiers.\\
A15 & Comparator fallback has a proof-obligation gap
& Audit candidate & S, P, U. No wrong public result established; strengthen the invariant.\\
A16 & Repeated symbolic expansion and linear-power paths deserve profiling
& Medium & S, P, U. Optimize only after contract-matched measurements.\\
\bottomrule
\end{longtable}

There is deliberately no entry claiming that the library cannot handle
fractional powers, logarithmic inverses, special functions, flat sectors,
or automatic series arithmetic. Those capabilities are already present.
There is also no claim that all built-in Wolfram asymptotic results carry a
rigorous numeric error bound. Comparisons must distinguish representation,
formal expansion, asymptotic theorem, and numerical certification.

\section{Detailed defects and reproducible contract examples}
\label{sec:details}

\finding{A01}{A sparse answer triggers a dense allocation}{High-priority resource defect; S, D, E; native reproduction U.}
\textbf{Location.} In the core module, \code{makeForwardObject} calls
\code{makeSeriesData} while constructing every ordinary forward result.
\code{makeInverseSeriesData} performs the analogous conversion for eligible
inverse results, including results reconstructed by refinement. Both form a
common rational-exponent denominator and then allocate a coefficient list.
The essential operation is:
\begin{lstlisting}
den = LCM @@ (Denominator /@ Append[exps, remData[[1]]]);
nmax = remData[[1]] den;
coeffs = Table[0, {nmax - nmin}];
\end{lstlisting}
No dense-span budget is checked at this point.
\cite{core,refinement}

\textbf{A minimal family.} Let $q\geq2$ be an integer and consider
\begin{lstlisting}
AsymptoticExpansion[1 + x^(1/q) + x, {x, 0, 1},
  "MaxTerms" -> 10]
\end{lstlisting}
The exclusive cutoff retains just $1$ and $x^{1/q}$; the first omitted
power is $1$. Thus the exported lattice has denominator $q$, lower index
zero, and upper index $q$. The requested coefficient-list length is
\[
 n_{\max}-n_{\min}=q.
\]
For $q=10^9$, this is a billion slots, despite only two retained sparse
blocks. The input is exact and compact, and its mathematically correct
sparse answer is trivial. This observation does not depend on estimates of
Wolfram's per-slot memory layout.

The independent checker computes these spans without allocating the lists.
The supplied native regression uses the much smaller denominator $10^6$
and wraps evaluation in time and memory limits. The billion-slot case is
\emph{not} executed, and no empirical memory number is claimed.

\textbf{Why this is more than a slow optional conversion.} The conversion
is eager: a user who never requests \code{s["SeriesData"]} can encounter
its cost. Moreover, \code{"MaxTerms"} bounds several sparse operations but
not this expansion of the representation. A resource guard in a preceding
sparse product does not protect the constructor from a later dense export.

\textbf{Recommended correction.} Preserve the sparse answer and make
classical conversion an explicit, lazy operation. Compute the common
lattice and its span first; refuse conversion when the representation
budget is exceeded. Return an informative \code{Missing} or \code{Failure}
from the export request, not from the otherwise successful sparse
construction. Cache a successful conversion only on explicit demand, with
a versioned cache policy. The plan should report retained blocks, dense
slots, denominator, and whether an error contract would be lost.

The supplied patch generator adds a conservative 200,000-slot guard to
both exporters. This is an immediate containment proposal, not a substitute
for a configurable lazy-export design. It is deliberately limited to the
pinned source and does not modify the repository in this review.

\textbf{Acceptance criteria.} Construction of the sparse family must remain
proportional to its sparse representation, subject to exponent arithmetic;
an optional conversion above its budget must return a structured reason;
existing small classical exports must remain unchanged; inverse construction
and refinement must obey the same rule. Tests must run against both modular
and regenerated standalone entry points.

\finding{A02}{Logarithmic remainder information disappears on export}{High-priority interoperability-contract defect; S, D, E; native reproduction U.}
The rich remainder descriptor contains both a power and a logarithmic
degree. The two exporters use the power but not
\code{remData[[2]]}. Therefore they cannot preserve that two-dimensional
magnitude contract in their current return value.
\cite{core}

Consider the finite source expression and cutoff
\begin{lstlisting}
s = AsymptoticExpansion[1 + x Log[x], {x, 0, 1}];
{s["RemainderPower"], s["RemainderLogDegree"]}
(* Expected rich metadata: {1, 1}. Native test supplied. *)
\end{lstlisting}
The discarded term is $x\log x$, so the rich object's natural bound is
$\OO(x(1+|\log x|))$. The exporter constructs the same ordinary coefficient
list and power frontier that it would construct after discarding $x$ from
$1+x$. The missing distinction is mathematically material:
\[
 \frac{|x\log x|}{x}=|\log x|\longrightarrow\infty,
 \qquad x\downarrow0.
\]
Thus an algebraic magnitude bound $\OO(x)$ cannot replace the richer one.

\textbf{Important qualification.} This is not an accusation that native
\code{SeriesData} is intrinsically incorrect for logarithmic expansions.
Its formal series-order conventions and the package's explicit magnitude
contract are different interfaces. The defect is that a conversion
advertised as a useful property silently omits a piece of the package's
contract, and a downstream consumer cannot reconstruct it from the
exported object alone. Unlike \code{Normal}, the export does not make the
loss explicit.

Three policies are defensible. A conservative export returns
\code{Missing["LogarithmicRemainder"]} whenever the logarithmic degree
cannot be represented faithfully. A richer export can return both the
classical object and a remainder descriptor with an explicit compatibility
statement. A third, deliberately weaker, export can replace the power by
$\rho-\varepsilon$ for an explicitly chosen $\varepsilon>0$, since
\[
 u^\rho(1+|\log u|)^d=\OO(u^{\rho-\varepsilon}).
\]
The third choice must recheck the retained support: the weakened frontier
cannot be inserted blindly below terms that remain in the classical list.

The supplied containment patch chooses the first policy. Tests should cover
zero and nonzero logarithmic degrees, empty retained lists, inverse exports,
and results whose coefficients themselves contain logarithms. A coefficient
containing a logarithm is not by itself a reason to reject a formal series;
the specific concern is the \emph{remainder contract}.

\finding{A03}{Explicit observable powers are silently shadowed}{Medium-priority option-contract ambiguity; S, D; native reproduction U.}
\code{InverseExpansionCoefficient} accepts either a normalized model or an
inverse object. Its model overload reads the public \code{"Power"}
option. Its object overload delegates as follows:
\begin{lstlisting}
InverseExpansionCoefficient[a["Model"], k,
  "Power" -> inheritedInternalPower, opts]
\end{lstlisting}
An explicit caller option is placed after an already supplied value.
Wolfram's option-precedence behavior makes that inherited value effective.
\cite{core,optionvalue}

For $f(x)=x+x^2$, let $g=f^{-1}$ near zero. Direct algebra gives
\[
 g(y)=y-y^2+2y^3+\cdots,
 \qquad g(y)^2=y^2-2y^3+\cdots.
\]
The first correction coefficient is $-1$ for power one and $-2$ for power
two. Consequently, an object query that accepts \code{"Power" -> 2}
but retains the original power-one computation is observably different
from the corresponding model query. The independent checker verifies the
algebra; the native suite supplies both queries.

There are two reasonable API contracts. An object coefficient query can
mean ``a coefficient of exactly this stored observable'', in which case a
conflicting power option should be rejected. Alternatively, it can mean
``a coefficient from this stored forward model for the requested
observable'', in which case an explicit power should override inheritance.
Neither contract justifies silently accepting an ineffective option.

The included helper \code{InverseCoefficientForPower} implements the
second policy for ordinary power--log inverse objects. It first resolves
the \emph{external} requested observable power and only then converts to
the internal exponent required at a source infinity. Simply moving the
option to the front is not an adequate general repair: the sign conversion
must still be applied to the explicit value. The helper also rejects
unsupported kinds instead of indexing a model they do not carry.

\textbf{Acceptance criteria.} Test inherited and explicit powers at finite
and infinite source endpoints, negative branches, invalid powers, model
queries, and nonordinary result kinds. Document whether a coefficient query
changes the observable or only inspects an existing object. Any chosen
policy should be consistent across the general and Fourier-specific
coefficient interfaces.

\finding{A04}{Refinement can return a valid result below its requested goal}{Medium-priority goal-satisfaction defect; S, D, E; native reproduction U.}
The generic operation-replay branch of \code{SeriesRefine} refines each
operand using a guard margin based on the requested cutoff and the old
operand precision. It does not generally solve the inverse precision
problem for the operation. The source uses a margin of the form
\[
 h+2+|\min(0,\rho_{\rm old})|.
\]
The arithmetic afterward still transports the available remainders. This
is good for correctness of the returned bound, but does not ensure that the
requested goal has been achieved.
\cite{seriesops}

A concrete source trace is:
\begin{lstlisting}
s = AsymptoticExpansion[Exp[x], {x, 0, 2}];
t = AsymptoticExpansion[x^-100, {x, 0, 1}];
p = SeriesMultiply[s, t];
r = SeriesRefine[p, 5];
\end{lstlisting}
The exact pole operand requires no unknown-tail improvement. For the
exponential operand, the generic replay asks for cutoff $5+2=7$. Multiplying
its $\OO(x^7)$ tail by $x^{-100}$ gives only $\OO(x^{-93})$, whereas the
requested product goal is $\OO(x^5)$. The exponential operand actually needs
precision $105$.

This is \emph{not} a claim that the multiplication invents coefficients or
returns an invalid bound. Its precision calculus clips to what is known.
The problem is a missing guarantee or explicit failure at the refinement
request boundary. An object can carry the requested goal in its history
without achieving that goal in its returned error grade.

The right repair is backward precision planning. For a product with leading
valuations $\alpha=\nu(A)$ and $\beta=\nu(B)$, the first necessary power
requirements for output precision $H$ are
\[
 \rho_A\geq H-\beta,\qquad \rho_B\geq H-\alpha.
\]
One must also check the error--error product and logarithmic envelope
degrees. If the available leading valuation is unknown because of
cancellation, obtain a leading-term certificate or refine it first; do not
silently assume zero valuation. The supplied Python check and Wolfram
helper verify the simple $105$ demand, but do not claim to implement the
complete planner.

\textbf{Acceptance criteria.} A refinement request should return either a
result whose reported achieved grade meets the requested grade, or an
explicit failure carrying the best available result and the reason for
shortfall. Tests should include large poles, reciprocals near cancellation,
powers of a vanishing series, compositions with large valuation multipliers,
and operations mixing exact carriers with asymptotic brackets.

\finding{A05}{Normalization hides an integer-factorization problem}{Medium-priority performance and predictability risk; S, P; native timing U.}
The core's \code{logCanon} expands logarithms of positive integer and
rational constants using \code{FactorInteger}. It is called within
coefficient and polynomial canonicalization. Thus an input containing an
otherwise harmless constant such as $\log N$ can introduce the cost of
factoring $N$ before the requested asymptotic computation needs that
factorization.
\cite{core}

For a family $N=pq$ with large distinct primes, the expression $\log N$ has
a compact exact representation. Factoring it is unnecessary for returning
an expansion such as $1+(\log N)x+\OO(x^2)$. The concern is the coupling of
routine normalization to an expensive unrelated algebraic task, not a
measured assertion that every large integer is difficult. No expensive
factorization fixture was run in this audit.

A safer normalizer keeps $\log N$ atomic by default, performs cheap rational
and structural simplifications first, and invokes factorization only under
an explicit effort budget or for sufficiently small constants. A
request-local cache can avoid repeating successful work. A cache of failed
or timed-out canonicalizations must include the assumptions and effort
policy so that a later stronger request is not suppressed accidentally.

This change creates a genuine tradeoff: aggressive normalization can prove
that two independently written coefficients cancel, and therefore can
improve the selected frontier. The solution is not to replace exact equality
with numerical equality. It is to make unresolved equality and expensive
normalization visible, preserve a conservative remainder when a cancellation
is unproved, and allow users to opt into stronger simplification.

\section{Audit of mathematical and implementation invariants}
\label{sec:invariants}
\subsection{Sparse ordering, resonances, and the comparator candidate}
The sparse engine merges equal weights, combines coefficient polynomials,
and sorts the result. Algebraic canonicalization is important: for example,
$\sqrt8$ and $2\sqrt2$ must not create different layers. Multi-index
resonances are equally important. If gaps are $1$ and $2$, the multi-indices
$(2,0)$ and $(0,1)$ contribute to the same weight and must be combined before
the layer is counted, declared zero, or omitted.

The implementation already contains complete-layer and complete-boundary
logic, including retained refinement state. These are substantive features,
not proposed additions. Testing should stress their invariance under
permutation of forward blocks, algebraically equivalent weights, and
cancellation of an entire first omitted layer.
\cite{core,refinement,fourier}

A narrowly scoped candidate deserves attention in the exact comparator.
After several earlier equality and sign tests, the fallback simplifies
$d<0$ and returns a positive comparison when that predicate simplifies to
\code{False}. Logically, $\neg(d<0)$ proves only $d\geq0$, not $d>0$.
The earlier equality tests may already settle this for the intended exact
numeric inputs; this review has not established a public counterexample.
The finding is therefore A15, an invariant-hardening opportunity, not a
confirmed incorrect exponent order.
\cite{core}

The robust contract is a three-way proof procedure: establish $d=0$,
establish $d<0$, or establish $d>0$. Otherwise return an undecidable-order
failure. For symbolic weights under assumptions, keep that procedure
separate from a comparator whose precondition is a fixed exact real number.
Never repair the gap using a fixed-precision numerical sign test near zero.

\subsection{Native-series ingestion is an analytic boundary}
The native special-function adapter is not merely a call to
\code{Normal[Series[...]]}. It recognizes candidate heads, obtains native
series information, and then applies checks on the supported representation,
realness, coefficients, and error envelopes. A list of recognized heads is
not a promise of all parameter values and endpoints for those functions.
The native adapter also contains a conservative weakening of a power
frontier to absorb logarithmic growth in a tail.
\cite{native}

That boundary needs especially adversarial tests. A native expression can
contain a nested or multiplied \code{SeriesData}, a retained exponential
carrier, a branch-dependent coefficient, or a coefficient whose apparent
zero is conditional on parameters. Every accepted conversion should record
which part of the native result justified the retained blocks and which
part justified the error envelope. A fallback must not turn a missing
native remainder into a guessed exact finite function.

Useful cases include expansions with their first several coefficients
zero, logarithmic singularities at a finite point, parameter values that
cancel the generic leading coefficient, and special-function identities
that are real only on a restricted source approach. The review does not
establish a wrong native-ingress result for such a case; these are
high-value regression candidates. A deliberately unsupported case returning
a precise failure is preferable to an unjustified broad acceptance claim.

\subsection{Fourier coefficients and oscillatory zeros}
The Fourier engine stores finite Fourier-polynomial modes in the logarithmic
variable and validates real conjugacy relations. It refuses an oscillatory
leading block without an eventual nonzero sign. Its remainder explanation
uses coefficient norms and polynomial envelopes, not the value of one
oscillatory factor at a numerical sample point. This is the appropriate
principle: a vanishing $\sin(\omega\log u)$ is not evidence that an entire
unknown remainder vanishes.
\cite{fourier}

Future Fourier work should preserve exact frequencies through products and
compositions. Frequency truncation needs its own error bound; a
\code{"MaxFrequencies"} resource failure must not be replaced by dropping
small-looking modes. Test zero-frequency cancellation, conjugate pairs,
repeated frequencies, resonant source weights, and agreement of the general
residual dispatcher with \code{FourierInverseResidual} on supported objects.
The dedicated Fourier coefficient accessor already exists and should be
advertised by a capability query rather than obscured behind a generic
wrong-kind failure.

\subsection{Finite flat sectors are not a universal transseries field}
The repository has explicit finite flat-sector construction and operations.
That is a significant extension beyond an ordinary power--log jet, but it
should remain described as its implemented finite scale calculus. A general
well-based logarithmic--exponential transseries field requires closure,
comparison, composition, and support-management guarantees beyond a set of
family-specific reductions. Those guarantees cannot be inferred from the
presence of a function named for a transseries-related operation.
\cite{readme,seriesarith}

A good invariant is that omitted sectors are represented by a genuine
sector envelope, not by an algebraic power that cannot dominate the omitted
exponential on the selected approach. Cancellation of exact sectors should
not erase an uncertain empty jet. The current arithmetic code has distinct
flat and composite routes; changes should preserve that distinction instead
of forcing all results into an inappropriate single power cutoff.

\subsection{What the exact interval certificate proves}
The certificate evaluator encloses elementary expressions with exact
rational endpoints and directed dyadic rounding. Exponential and logarithmic
enclosures use explicit series tails; products, reciprocals, and permitted
powers preserve interval inclusion. The code checks the selected source
side and the retained source-domain condition on the entire closed interval.
It then requires the derivative to be separated from zero.
\cite{certificates}

The underlying argument can be stated independently. Suppose $f$ is
continuous on an interval $I$ and differentiable in its interior, with
$|f'|\geq\mu>0$ and a constant derivative sign. For a center $c\in I$, an
exact residual enclosure implies $|f(c)-y|\leq\varepsilon$. If
$[c-\varepsilon/\mu,c+\varepsilon/\mu]\subseteq I$, monotonicity and the
mean-value theorem establish a root between those endpoints, and strict
monotonicity gives uniqueness. Alternatively, certified endpoint signs can
establish existence. A signed derivative enclosure can then sharpen the
root interval. These are the two existence routes visible in the code.

For example, for $f(x)=x+x^2$, $y=1/10$, $c=1/10$, and $I=[0,1/5]$, the
residual is $1/100$ and $f'\geq1$. The bracket $[9/100,11/100]$ lies in $I$.
The independent check verifies these rational inequalities; it does not
execute the package certificate function.

Crucially, the returned certificate says that it does not certify the whole
family of functions compatible with a declared \code{InputRemainder}.
Nor does uniqueness on one verification interval establish uniqueness of a
global inverse branch. These limitations are already made explicit by the
implementation and should remain prominent in the user interface.
\cite{certificates}

An independent certificate checker would improve assurance further. The
current evaluator is mathematical software, not a machine-checked proof of
its own implementation. A compact transcript containing the exact interval
operations, tail inequalities, domain checks, and final containment test
could be replayed by a smaller checker. Such a checker would validate the
particular certificate, not establish correctness of every heuristic used
to choose its initial center.

\subsection{Distinguishing three kinds of residual evidence}
A formal jet residual verifies algebraic cancellation below a specified
grade in a stored finite model. A numerical residual or a numerical
reference-root comparison supplies finite-precision evidence at selected
points. An interval certificate proves a local statement for an explicit
function with the enclosures it actually supports. The library exposes all
three kinds, and their names, return fields, and documentation should keep
them separate.
\cite{core,certificates,special}

A zero residual in a finite Stirling model does not remove the original
Stirling tail. A tiny residual need not imply a small inverse error when
$f'$ is small. Agreement at a handful of points cannot resolve a branch
ambiguity. These are reasons to improve provenance presentation, not reasons
to discard the useful formal and numerical tools.

\section{Performance analysis and resource design}
\label{sec:performance}
\subsection{The relevant complexity parameters}
A benchmark indexed only by the requested number of terms misses most of
the important costs. At least six quantities matter: retained support size,
number of candidate products, logarithmic polynomial degree, arithmetic
complexity of exponents, number of independent gaps, and complexity of the
symbolic coefficient field. Classical export adds a seventh, the dense
lattice span.

For $m$ positive gaps, the multi-index set
\[
 \{k\in\N^m:k\cdot\delta<H\}
\]
has a leading volume estimate
$H^m/(m!\prod_j\delta_j)$ as $H$ grows with fixed positive gaps. This is a
lattice-point estimate, not a package timing prediction. Small gaps or many
independent gaps can create many candidate indices even when resonances
merge them into comparatively few displayed powers. An implementation that
limits only returned blocks does not bound this work.

For jets of sizes $n$ and $m$, an untruncated product has $nm$ candidate
pairs before merging. The current \code{jetMul} already uses sorted weights
and a monotone scan to skip products beyond the cutoff and to check a pair
budget before multiplying the coefficient polynomials. Similar pruning was
added to Fourier products. These optimizations should be retained and
regression-tested; proposing a Cartesian product followed by filtering
would move backward.
\cite{core,fourier,validation}

\subsection{A06: one option is not one budget}
The public \code{"MaxTerms"} value is used in different places for retained
or candidate terms, index regions, product counts, polynomial expression
size, and model orders. Other paths use a frequency cap, hard-coded
simplification timeouts, or chart-recursion limits. A user who raises
\code{"MaxTerms"} cannot predict which resource has actually been granted.
\cite{core,native,gamma,sourcecoords,refinement}

A structured budget should distinguish at least
\code{MaxBlocks}, \code{MaxCandidateProducts}, \code{MaxIndices},
\code{MaxCoefficientLeaves}, \code{MaxDenseCoefficients},
\code{MaxFrequencies}, and a total time or work budget. The existing option
can remain a backwards-compatible shorthand during migration. Every
resource failure should name the dimension, the requested or observed
amount where available, and whether any validated partial result exists.

Local symbolic subcalls should consume the caller's remaining budget
rather than each receiving a fresh hidden allowance. A timeout should mean
``proof or computation not completed under this budget'', not a negative
mathematical conclusion. In particular, a timed-out sign proof cannot be
used as evidence of nonpositivity, and a failed simplification cannot be
used as evidence that two expressions differ.

\subsection{A16: candidates for improvement, not measured speedups}
The native adapter includes a repeated-multiplication path for positive
integer powers, while the ordinary and Fourier algebras have more efficient
power-building structures. Binary powering with early truncation is a
reasonable candidate for the native path, but its envelope normalization
cost must also be measured. Intermediate expression growth can dominate the
number of multiplications.
\cite{native,core,fourier}

The Gamma inverse coefficient builder repeatedly constructs and expands a
phase through successively increasing orders. This is easy to understand
and valuable as a reference implementation. A retained coefficient cache,
incremental truncated polynomial arithmetic, or a Newton-style update could
reduce repeated work. However, the main recommendation is to benchmark
representative orders before replacing a transparent recurrence. The
current code's \code{LeafCount} guard is useful but does not necessarily
bound the cost of a symbolic expansion that occurs before the guard.
\cite{gamma}

Coefficient normalization should be staged. Cheap exact arithmetic,
structural zero checks, and algebraic canonicalization should precede
expensive global simplification or factorization. Request-local memoization
can reuse exponent comparisons and coefficient normal forms. Cache keys
must include assumptions and the relevant normalization policy. Do not
introduce an unbounded global cache of expression trees or silently reuse
facts established under incompatible assumptions.

For ordinary inverse refinement, model signatures, retained coefficient
prefixes, and residual checks of recovered state already exist. This is a
stronger starting point than replaying every request from zero. An automatic
method selector can build on \code{Lagrange}, \code{Newton}, and grouped
methods already present: estimate index-region size, degree growth, and
expected resonance before selecting a route. The selected method and any
fallback should be recorded in a plan.
\cite{refinement,core}

\subsection{Benchmark protocol and the supplied harness}
The included \code{CompareBuiltins.wls} measures selected package and native
operations, subject to per-evaluation time and memory limits. It records
actual outputs, messages, result heads, object size, finite-expression size,
and time for extracting the normal expression. Its ``operation-cold''
measurement clears the symbolic system cache; it is explicitly \emph{not}
a fresh-process cold start. Warm observations are recorded separately.
No numbers from that harness are present because it has not been run here.

A credible future comparison should separate loading, construction,
normalization, presentation, classical export, and refinement. The package
object contains metadata and sometimes reusable state, whereas a native
result may be just a formula. Comparing only their total byte counts or
construction times can be misleading. Compare normal expressions and
achieved mathematical information as well as complete objects.

The two systems' order arguments are not automatically equivalent. An
exclusive package cutoff of four does not necessarily match a native
``order four'' request. Select a common retained prefix or prove residual
agreement at a common grade before computing a speed ratio. For implicit
native solutions, branch selection and the number of returned branches are
part of that matching task.

Suggested workloads are small Taylor/Puiseux baselines; irrational gaps;
resonant logarithmic inverses; high-denominator sparse powers; large-pole
refinement; cancellation-heavy arithmetic; Fourier frequency growth; and
Gamma/Barnes exact-core inverses. Use exact data, fixed assumptions, pinned
kernel and package versions, multiple repetitions, and recorded failures.
A timeout is a censored observation, not an arbitrary large successful
runtime. Repository-reported operation improvements should remain labeled
as local package-to-package measurements, not generalized into built-in
comparisons.
\cite{validation}

\section{API and user-experience redesign}
\label{sec:ux}
\subsection{A09: make the order request a typed mathematical object}
The same-looking cutoff argument has different meanings in different
families. That flexibility is mathematically natural, but it is difficult
to discover from a call that selects its scale automatically.
\cite{core,special,gamma,readme}

\begin{longtable}{@{}P{0.26\linewidth}P{0.31\linewidth}P{0.36\linewidth}@{}}
\toprule
Representation & Order coordinate & Important qualification\\\midrule
\endhead
Ordinary power--log & Positive local source or target power & Cutoff is
exclusive; coefficients are complete logarithmic blocks.\\
Lambert/logarithmic reduction & Inverse-logarithmic variable in a unit
bracket & The extracted prefactor is not counted as an ordinary power of
the original target.\\
Gamma/Barnes ordered inverse & Reciprocal exact core & Each core-power
block retains its polynomial in an inverse logarithm.\\
Legacy explicit Gamma special adapter & Exact-core perturbation marker &
The integer argument denotes an inclusive marker depth, not the same
ordered core-power cutoff.\\
Finite flat sectors & Sector support and within-sector grades & A single
ordinary power cutoff cannot describe all omitted sectors.\\
Composite error envelope & Recorded compatible approach & There may be no
single scalar cutoff; operand truncation must be explicit.\\
\bottomrule
\end{longtable}

Retain concise calls for standard use, but expose the resolved request:
\begin{lstlisting}
(* Proposed interface, not an implemented repository feature. *)
SeriesPlan[f, {x, x0},
  <|"Goal" -> "PowerCutoff", "Value" -> h,
    "Coordinate" -> Automatic, "Boundary" -> "Exclusive"|>]
\end{lstlisting}
A plan should explain the source approach, target approach, selected scale,
exact carrier, meaning of a block, transported input ceiling, expected
backend, and resource limits. A result should contain both
\code{RequestedPrecision} and \code{AchievedPrecision}, together with an
explicit \code{GoalSatisfied} value. A request that cannot be meaningfully
compared to the selected scale should fail before performing substantial
coefficient work.

The project already accepts structured refinement requests for additional
ordinary inverse blocks and for certified numerical tolerances. The
additional-block route reports whether the goal was reached, preserves a
best expansion on resource failure, and distinguishes exact termination.
The numerical route explicitly says that its certificate does not improve
the symbolic series remainder. The proposed general precision interface
should extend this existing design to numeric-cutoff operation replay,
not replace it or claim to introduce structured requests for the first time.
\cite{refinementrequests}

A term goal also needs a cancellation policy. ``First $n$ nonzero complete
blocks'' is not the same as ``first $n$ multi-indices'' or ``first $n$
monomials after expanding every logarithmic polynomial''. The package's
complete-block convention is a good default; display it alongside the
returned count. For exact finite answers, returning fewer blocks should
report exact termination rather than look like an incomplete computation.

\subsection{A10: keep approximate evaluation distinct from checked evaluation}
The object accessor \code{s[value]} substitutes into the finite expression.
It does not call a domain prover, solve the original equation, or return an
error enclosure. This is documented behavior, so it is not an implementation
bug. Nevertheless, the notebook presentation can make it look more
certified than it is, especially when the head and metadata are hidden.
\cite{core}

A separate proposed \code{SeriesEvaluate} interface should distinguish a
finite approximation, a numerical diagnostic, and a certified enclosure.
An approximation result should report whether the point satisfies known
domain predicates, whether that check is proved or merely numerical, and
whether an error bound is quantitative or only asymptotic. A symbolic
Big-O class without a constant and validity radius cannot produce certified
error bars at a chosen point.

The inexpensive accessor can remain, but a compact summary panel should
make its status clear: scale, branch, retained blocks, known precision,
certificate availability, and whether the displayed formula has dropped
anything. Do not place an unconditional green ``certified'' badge on a
series whose only evidence is a formal zero residual.

\subsection{A11: version the result schema and expose capabilities}
Different result kinds carry different fields. An ordinary inverse has a
\code{Model}; a Fourier inverse has frequency-specific normalized data;
a special-function object can retain a coordinate series and an exact
carrier. Operations currently dispatch by combinations of kind, scale,
and field presence. This works, but makes maintenance and error messages
fragile as new families are added.
\cite{core,seriesops,fourier,special}

Introduce a small validated schema with stable mandatory fields and
kind-specific extensions. A capability query should answer whether an
object supports ordered truncation, additional blocks, formal residuals,
coefficient queries, classical export, differentiation, and numeric
certification. Each answer should carry a reason and, where appropriate,
the additional required hypotheses. Unsupported operations should return a
precise kind-aware failure, not a generic invalid-argument message arising
from delegation to a missing model.

The 1.8.0 validation record states that \code{GeneralizedSeries} replaced
the previous public head and that the old head is not exported as an alias.
That intentionally affects explicit patterns and saved input. Provide a
migration function for supported saved structures and document which older
representations cannot be safely upgraded. Silent rewriting of arbitrary
associations would be dangerous: changed field meanings require validation,
not just head replacement.
\cite{validation}

Version computation state separately from display data. Serialized state
should include the package/schema version, model signature, assumptions,
branch, coordinate definitions, normalization policy, and source precision.
The current signature and compatibility checks provide a useful foundation.
User editing of metadata should invalidate cached state unless a validator
reestablishes the required invariants.
\cite{refinement}

\subsection{A12: normalize native conventions without hiding differences}
The ordinary local-coordinate function accepts \code{Automatic},
\code{"FromAbove"}, and \code{"FromBelow"}, rejecting other direction
forms. Numeric Wolfram direction aliases are not accepted by that path.
This is an interoperability omission, not a discovered wrong-branch answer.
A central normalization layer can map supported aliases to the package's
real approach representation before dispatch, while continuing to reject
complex directions outside its contract.
\cite{core,asymptoticsolve}

Assumption handling should also be explicit. Several public option lists
use \code{Assumptions -> True}, whereas native functions commonly use
\code{$Assumptions}. Some internal simplification calls may still interact
with the global environment; this review does not assert exact behavior
under every \code{Assuming} wrapper without native execution. Add tests
that compare explicit assumptions, a dynamic assumptions block, and a
retained object's later refinement. Choose a policy and retain the effective
assumptions with the result.
\cite{core,asymptotic}

Use consistent failure categories for malformed syntax, unsupported scale,
unproved branch condition, insufficient input precision, exhausted budget,
and numerical diagnostic failure. These conditions suggest different
remedies. For example, increasing the number of terms cannot repair a
missing derivative hypothesis, and increasing a time budget cannot justify
a complex branch in a real-only scale.

\subsection{A13--A14: certificate discovery and distribution polish}
\code{InverseCertificate} has an \code{"Interval" -> Automatic} option
default, but its current implementation requires an explicitly supplied
ordered rational interval. The usage text does show that interval, so this
is a low-priority sentinel-design issue, not a hidden undocumented
mathematical requirement. A missing-interval diagnostic should say what is
required. A future automatic bracket search must preserve the selected
branch and prove every successful bracket; failure to find one is not proof
that no inverse exists.
\cite{certificates,core}

The paclet manifest declares only a kernel extension. The repository's
Markdown/HTML guide is substantial, and the validation record reports that
all 37 public symbols are documented. Native help-system integration,
searchable symbol pages, executable examples, and a guide organized by
scale and supported operation would improve discoverability without
replacing those materials.
\cite{paclet,validation}

The root license text is titled ``MIT No Attribution License'', while the
paclet metadata says ``MIT''. Align these literal metadata identifiers so
packaging tools and users see a consistent declaration. This is a
metadata-consistency observation, not a legal opinion about licensing
obligations.
\cite{license,paclet}

The project already documents problems with direct cold HTTP loading of a
large standalone file and recommends downloading completely before
\code{Get}. Preserve that mitigation. Release instructions should offer a
commit-pinned file or paclet with a published checksum and an offline path,
rather than encouraging evaluation of a moving \code{main} URL. The
standalone builder's dependency checks and embedded module hashes are
valuable and should remain part of the release process.
\cite{development}

\section{Testing, CI, and release evidence}
\label{sec:testing}
\subsection{A07: packaging CI is not mathematical CI}
The inspected \code{standalone.yml} workflow checks that the generated
standalone source is fresh and runs the Python builder tests. These checks
are useful: stale generated code is a real distribution defect. They do
not execute the Wolfram package or compare its mathematical outputs.
The latest focused native validation is a separate, repository-reported
activity, not a native gate in this workflow.
\cite{ci,validation}

Add a native job on a properly licensed runner, or a release procedure with
equivalent auditable evidence when such a runner is not available for every
pull request. The release gate should load a clean pinned checkout,
regenerate and compare the standalone file, run both focused and complete
native suites, exercise the modular and standalone public entry points,
and retain the actual kernel version, platform, source hashes, selected
tests, counts, messages, and failures. Unsupported or skipped tests must
remain visible.

A useful sequence is fast structural and builder checks on every relevant
change, a focused native matrix selected from touched subsystems, and a
full clean-process native suite before release. A focused pass must not be
relabeled as a full pass merely because all tests selected by that runner
succeeded. The current validation documentation is commendably explicit
about this scope; automation should preserve the same honesty.

\subsection{A08: bring the full-suite runner up to the focused-runner standard}
The inspected \code{Tests/RunTests.wl} enumerates \code{*.wlt} files and
exits according to the failed-test count. It does not first reject an empty
file list or require a positive executed-test count. Therefore an accidental
empty selection is not explicitly prevented from appearing successful by
this runner. It also lacks the source-hash attestation described for the
newer focused-runner infrastructure.
\cite{runtests,validation}

The repository already has a stronger focused runner that checks loading,
rejects empty suites, records sources before and after execution, and
propagates export failures. Reuse or extract that infrastructure instead of
creating yet another independent runner. The finding does not apply to all
upstream test scripts. Add a fixture that deliberately points to an empty
test directory and requires a failing process status; also test missing
output directories and report-export failures.

\subsection{The highest-value new test classes}
A coefficient oracle alone is insufficient for a library whose main value
is metadata. Every representative fixture should check the expression,
selected approach, coordinate, remainder scale, exactness status, achieved
precision, and provenance appropriate to its kind. A correct displayed
formula with a false or lost contract is a failing result.

Property-based tests should randomize small exact finite models and verify
that the inverse residual vanishes below the requested internal grade.
Compare independent Lagrange and Newton routes on their overlap, then
compare a classical specialization with native reversion. Random generation
must preserve admissibility; an expected failure on a genuinely unsupported
scale is not a failed mathematical algorithm.

Metamorphic tests are especially useful here. Permuting or regrouping input
terms, replacing an algebraic exponent by an equal expression, applying a
known affine coordinate change, and refining in one step versus several
steps should preserve the appropriate contract. Truncation must never
strengthen an unknown remainder without new evidence. A refinement cache
with changed assumptions or a changed branch must be rejected. A generated
standalone package must agree with the canonical modular entry point.

Negative tests should cover invalid numeric directions, approximate input,
nonreal branches, ambiguous leading signs, missing derivative bounds,
input-precision ceilings, dense representation budgets, frequency budgets,
and timed-out proof attempts. Check failure tags and the absence of noisy
internal messages, not only whether the result matches \code{_Failure}.

Numerical tests should use exact starting parameters, adequate working
precision, and multiple approach points, ideally with a normalized residual
or error divided by the claimed scale. Near zeros of an oscillatory factor,
use the established envelope rather than that factor itself. A failed
\code{FindRoot} reference computation should be recorded separately from
a failed asymptotic approximation.

\subsection{What was actually run for this report}
\begin{longtable}{@{}P{0.40\linewidth}P{0.17\linewidth}P{0.36\linewidth}@{}}
\toprule
Evidence & Result & Scope\\\midrule
\endhead
Independent mathematical checker & 16 passed; 0 failed & Python/SymPy
identities, span calculations, and precision examples.\\
Python audit-tool unit tests & 7 passed; 0 failed & Patch transformation,
source-shape rejection, exact span calculation, and Git-blob hashing.\\
Package native regression tests in this review & Not executed & Native
connector failed; no local Wolfram kernel.\\
Proposed Wolfram helpers and patched exporters & Not executed natively &
Provided as reviewable proposals, not validated replacements.\\
Built-in comparison benchmarks in this review & Not executed & Harness
provided; no measured speedup claimed.\\
Latest upstream 1.8.0 focused record & 305 passed; 0 failed, as reported &
Selected 18 files; full suite explicitly not run at that milestone.\\
\bottomrule
\end{longtable}
The included evidence files preserve these categories. In particular, there
is no fabricated \code{native-audit-results.json} or built-in benchmark
result. Those files are created only if the supplied native runners are
actually executed.

\section{Proposed changes and use of the accompanying artifacts}
\label{sec:artifacts}
\subsection{A deliberately small containment patch}
\code{patches/make_export_patch.py} verifies the pinned core's Git blob
identifier before proposing a modification. Its default behavior prints a
unified diff and does not write the repository. An explicit
\code{--in-place} request creates a backup and atomically replaces only
the canonical core file. Source drift, a missing definition, or an already
applied patch is rejected. Line endings are normalized for the source-hash
comparison.

The patch does two things in each classical exporter: it returns an explicit
missing representation for a nonzero logarithmic remainder degree, and it
checks the dense coefficient span before allocating. It does not change
coefficient mathematics, implement lazy properties, fix refinement planning,
or settle the observable-option policy. The patch is intentionally small
enough to review independently of those larger changes.

The exact exporter excerpts used by the patch-tool unit tests are included
with provenance and the upstream license. They are \emph{not} a complete
package and must not be loaded as one. The whole repository was not copied
into this archive. The source-shape tests validate the transformation of
those excerpts; they are not native execution of the modified package.

A typical workflow, using a separate local checkout at the pinned revision,
is:
\begin{lstlisting}
# From the audit archive; replace the path with your pinned checkout.
python patches/make_export_patch.py \
  /path/to/Asymptotic/AsymptoticInverse/Kernel/AsymptoticInverse.wl

# Only after reviewing the diff, request an actual local change:
python patches/make_export_patch.py \
  /path/to/Asymptotic/AsymptoticInverse/Kernel/AsymptoticInverse.wl \
  --in-place

# From the upstream repository root, rebuild the generated distribution:
python validation/build_standalone.py
python validation/build_standalone.py --check
\end{lstlisting}
Native regression tests and comparison of modular and standalone loading
must follow. No remote repository write or pull request was performed.

\subsection{Helpers and native test organization}
\code{code/AuditSupport.wl} contains proposed, namespaced helpers. They do
not install overrides into the upstream package. \code{DenseSeriesPlan}
preflights a rational lattice and its remainder compatibility without
allocating the dense coefficients. \code{InverseCoefficientForPower}
resolves an explicit observable power before source-infinity conversion.
\code{RequiredProductPrecision} states the necessary valuation-based power
demands for a product, with its logarithmic and error-product limitations
explicitly recorded.

The native suite separates baseline behavior from desired contracts.
\code{Baseline.wlt} tests classical and logarithmic inverses, formal
residuals, the rich remainder, model coefficient queries, and a simple
explicit-function certificate. \code{Support.wlt} targets the supplied
helpers. \code{ProposedContracts.wlt} encodes the desired behavior for
A01--A04; it is \emph{not} claimed to pass the pinned unmodified package.
The A03 test chooses explicit-option precedence, while the article also
recognizes the alternative of rejecting a conflicting option.

Set \code{ASYMPTOTIC_PACKAGE} to an existing local modular core or standalone
file, then use the supplied runner. It requires a nonempty suite, records
the actual kernel and source file hash, and exports per-test details. Its
reference commit is descriptive, not proof that a caller's entire checkout
is clean. A production release attestation should hash all loaded modules.
\begin{lstlisting}
# POSIX shell example; the README also describes the environment settings.
export ASYMPTOTIC_PACKAGE=/path/to/Asymptotic/AsymptoticInverse.wl
wolframscript -file tests/RunAudit.wls

# Desired future contracts (not expected to pass unmodified upstream):
export ASYMPTOTIC_AUDIT_SUITE=Proposed
wolframscript -file tests/RunAudit.wls

# Separate, unexecuted comparison harness:
wolframscript -file benchmarks/CompareBuiltins.wls
\end{lstlisting}

\subsection{Independent reproduction and document build}
The Python checker can be rerun without a Wolfram kernel. SymPy is its only
nonstandard Python dependency. The Python unit tests also load that checker
and consequently need SymPy. The article uses standard LaTeX packages and
can be compiled without the upstream repository or network access once
those packages are installed.
\begin{lstlisting}
python code/independent_checks.py
python -m unittest discover -s tests -p 'test_python_tools.py' -v
cd article
pdflatex -interaction=nonstopmode -halt-on-error article.tex
pdflatex -interaction=nonstopmode -halt-on-error article.tex
\end{lstlisting}
These commands reproduce the independent checks and the document, not the
native package validation that remains outstanding.

\section{Development roadmap with acceptance gates}
\label{sec:roadmap}
\subsection{Stage 1: preserve existing information and bound representation costs}
First contain A01 and A02. Make dense export bounded, and prevent an export
from silently discarding a logarithmic error contract. Select and document
the A03 observable-option policy. Add native baseline and regression
fixtures before broad refactoring, and ensure the generated standalone
source is regenerated from the same canonical modules.

The completion gate is not merely ``the new tests pass''. Existing small
exports, inverse expressions, remainder fields, and formatting must remain
consistent on the native suite. Resource tests must confirm that an optional
representation failure does not destroy a valid sparse result. A release
record must identify the complete source set actually tested.

\subsection{Stage 2: make requests and results comparable}
Introduce a precision descriptor that identifies the coordinate, exact
carrier, power grade, logarithmic grade, sector information where relevant,
and hypothesis status. Add requested-versus-achieved reporting. Implement
backward demands first for ordinary sums, products, and powers, then for
logarithm, exponential, and composition. Keep safe fallback results, but
label an unmet goal explicitly.

The acceptance suite should include the large-pole example, exact zero,
cancellation that hides a leading valuation, and source remainder ceilings.
For every operation, write a precise rule for how input precision implies
output precision. New rules must not erase input uncertainty merely because
the retained polynomials cancel.

\subsection{Stage 3: separate planning, algebra, evidence, and presentation}
A maintainable architecture can be organized around four interfaces. A
planner selects a scale and computes precision demands. A scale algebra
implements exact operations and support ordering. An evidence layer records
formal residuals, analytic hypotheses, numerical diagnostics, and interval
proofs. A presentation layer displays expressions and handles intentionally
lossy exports.

This is not a recommendation to rewrite every module at once. Introduce
adapters around existing implementations and preserve them as reference
backends. Start with ordinary power--log objects, where the invariants are
clearest. A scale interface should define support comparison, block
normalization, truncation, operation precision, and available evidence.
A registry then replaces scattered field-presence tests with explicit
capabilities.

The completion gate is agreement with the existing backend on supported
cases, clearer failures on unsupported combinations, and no regression in
state-signature checks. Only after those invariants are stable should the
new interface become the foundation for further scale families.

\subsection{Stage 4: optimize where measurements identify bottlenecks}
Profile canonicalization, index enumeration, coefficient polynomial work,
repeated Gamma expansions, and classical export separately. Add bounded
request-local caches and improved powering where justified. Select inverse
methods from estimated work, with reproducible fallback decisions.

The gate is a contract-matched benchmark suite with raw observations,
messages, achieved grades, and both cold-process and warm-operation modes.
A smaller expression or faster first call is not a success if it loses a
branch condition, a remainder degree, or refinement provenance.

\subsection{Stage 5: expand analytic scope with explicit proof obligations}
Add new adapters only when their admissible domain, forward model,
remainder, derivative information, coordinate transformation, and inverse
conditioning are specified. Each adapter should provide an unsupported-case
boundary and tests at degenerate parameter values. A broad recognized-head
list without those contracts is less useful than a smaller reliable domain.

Prioritize extensions that reuse existing exact cores and error transport.
Gamma/Barnes interval evidence, additional threshold models, and controlled
source-chart composition are natural opportunities. Complex sectors and
multivariate systems require a separate architectural decision rather than
an accumulation of exceptional dispatch rules.

\section{Longer-term mathematical opportunities}
\label{sec:future}
\subsection{Certified special-function inverses}
The repository already provides Gamma/Barnes and Erfc asymptotic machinery
and quantitative forward-model information in explicit adapters. The next
step is not ``implement these inverses''; it is to connect admissible
forward-tail bounds and derivative bounds to a trustworthy original-function
inverse enclosure. That may require additional exact constants, special
function enclosures, and interval-domain checks beyond the current
elementary certificate evaluator.
\cite{special,gamma,certificates}

An adapter contract can return an explicit model $F_m$, a bound
$|F-F_m|\leq B_m$ on a specified interval, and derivative information
sufficient to separate $F'$ from zero. The residual bound becomes the sum
of the model-evaluation enclosure and the forward-model error bound, rather
than treating $F_m$ as the original function. Any unknown multiplicative
constant in a bare Big-O must be resolved before this becomes a quantitative
certificate.

\subsection{Uniformity in parameters and transition models}
Current exact parameter assumptions help establish realness and signs, but
an expansion valid for each fixed parameter need not be uniform when that
parameter approaches a degeneracy. For instance, in
\[
 f_\varepsilon(x)=\varepsilon x+x^2,
\]
the ordinary inverse for fixed $\varepsilon>0$ starts with $y/\varepsilon$,
whereas at $\varepsilon=0$ the positive branch is $\sqrt y$. The crossover
occurs when $y$ is comparable to $\varepsilon^2$. No number of additional
terms in the fixed-parameter expansion makes it uniform across that
transition without changing the scale or retaining a suitable exact core.

A future plan can report parameter-dependent dominant balances and suggest
an exact transition core such as
$(-\varepsilon+\sqrt{\varepsilon^2+4y})/2$. Uniformity should be a separate
proved property, with a parameter region and scaled target domain. This is
a promising direction because it uses the project's core/perturbation
architecture rather than fighting it.

\subsection{Complex sectors and branch geometry}
Extending a real germ to complex sectors requires more than permitting
complex constants through input validation. Powers and logarithms need
branch data; asymptotic dominance depends on the approach direction;
exponentials can change from decaying to growing; and a real interval
certificate no longer supplies the relevant existence theorem.

A suitable representation would retain a sector, branch cut or logarithm
determination, and admissible comparison relation. Start with a narrow
sectorial family whose remainders are already known, and keep it distinct
from the real engine. Stokes transitions and exponentially small terms
should be described only to the extent actually implemented. The current
finite-sector or real-flat machinery is not by itself a resurgent-analysis
system.

\subsection{Multivariate implicit inversion}
A system version would replace the scalar derivative lower bound with
Jacobian invertibility and a quantified conditioning bound. Formal
coefficient construction can use multivariate graded jets or an ordered
monomial set, but support order and precision become partial-order problems.
The native implicit-asymptotic solver is a useful source of candidate
solutions and differential tests, not an automatic certificate of a new
package representation.
\cite{asymptoticsolve}

A practical first target is a finite-dimensional real system near a regular
point with a fixed weighted grading and a nonsingular exact leading
Jacobian. Singular systems, several competing dominant balances, and
parameter-dependent rank changes should remain separate milestones.

\subsection{Machine-checkable local evidence}
The elementary interval certificate is particularly suitable for a small
independent checker because its arithmetic endpoints are rational. A
transcript could encode elementary expression nodes, input intervals,
directed rounding witnesses, elementary tail bounds, derivative exclusion,
and the final root-existence argument. Replay would validate a certificate
without trusting the asymptotic planner or numerical seed generator.

For formal series identities, a second small checker could verify sparse
coefficient arithmetic and the residual grade from a supplied model.
Keeping these checkers separate respects the distinction between a formal
identity, an analytic remainder theorem, and a numerical existence proof.
A future proof-assistant integration could certify those small kernels
before attempting the much larger symbolic-dispatch and simplification
system.

\section{Conclusion}
The reviewed package has a coherent and valuable direction: explicit real
scales, complete coefficient blocks, preserved remainder information,
branch-aware inverse computation, and inspectable evidence. Its mathematical
scope is already broader than its historical name suggests, and several
important safeguards and performance improvements are already implemented.

The most actionable findings are at boundaries. Eager classical export can
turn a tiny sparse answer into a huge allocation; that export also loses
logarithmic tail information. A coefficient wrapper silently shadows an
explicit option, and a generic refinement path does not always obtain the
precision its caller requested. These issues are narrow enough to address
without replacing the underlying theory.

The recommended strategy is to protect those boundaries, establish a
current native release gate, and turn precision, capabilities, and evidence
into stable interfaces. Native Wolfram functionality should remain both a
backend and an independent comparison baseline wherever contracts overlap.
Future progress should be judged not just by the number of recognizable
functions, but by how clearly the system states what it knows, what it has
assumed, what it has computed, and what it has actually proved.

\appendix
\section{Source coverage and reproducibility map}
\label{sec:coverage}
This appendix distinguishes direct source inspection from inventory and
family-level assessment. It does not claim an exhaustive line-by-line audit
of every repository file.

\begin{longtable}{@{}P{0.36\linewidth}P{0.24\linewidth}P{0.33\linewidth}@{}}
\toprule
Source or group & Inspection scope & Role in this review\\\midrule
\endhead
Core \code{AsymptoticInverse.wl} & Substantial selected ranges throughout
& Input/exponent handling; sparse jets; forward/inverse assembly; exports;
properties; residuals; coefficient API; module dispatch.\\
\code{SeriesOperations.wl} & Substantial selected ranges & Representation,
arithmetic precision, differentiation, and generic refinement replay.\\
\code{SeriesArithmetic.wl} & Lines 1--165 & Held normalization, regular
operands, automatic arithmetic, and fallback classification.\\
\code{RefinementState.wl} & Lines 1--145 & State signatures, prefix reuse,
Newton state checks, result assembly, and precision ceilings.\\
\code{RefinementRequests.wl} & Complete short module & Additional-block
and numerical-certificate requests; confirms numeric-cutoff replay remains
in the generic operation path.\\
\code{FourierCoefficients.wl} & Lines 1--310 & Finite modes, realness,
construction, coefficient envelopes, residuals, and coefficient queries.\\
\code{NativeSpecialFunctions.wl} & Lines 1--190 & Native-series ingress,
normalization, and integer-power implementation candidates.\\
\code{InverseCertificates.wl} & Lines 1--355, with the last retrieved range
partially truncated & Rational interval arithmetic, domains, existence and
uniqueness, options, and scope. Later adaptive-loop tail not fully audited.\\
\code{SpecialFunctionAdapters.wl} & Lines 1--150 & Erfc model bounds,
explicit Gamma marker adapter, and threshold families.\\
\code{SourceCoordinates.wl} & Lines 1--110 & Log/exponential charts, exact
reconstruction, domain handling, and explicit input-remainder limits.\\
\code{GammaInverse.wl} & Lines 1--105 & Family recognition, ordered core
coefficients, and constructor checks.\\
Other kernel modules & Complete filename inventory; dispatch and current
feature descriptions & Architectural and scope assessment, not an
exhaustive implementation audit.\\
README, PacletInfo, license, CI, test runner, development notes & Directly
retrieved text & Public claims, packaging, migration, loading, and gates.\\
Current validation summary and historical review manifest & Retrieved
records, not native reruns & Scope and provenance of reported tests.\\
\bottomrule
\end{longtable}

The archive includes the full list of 38 kernel \code{.wl} module names and
\code{init.m}. Its source-coverage JSON records the pinned revision and the
retrieved ranges. Repository citations below link to that revision; official
Wolfram references describe native capabilities rather than an independently
measured installed kernel.

\section{Artifact inventory and status}
\begin{longtable}{@{}P{0.45\linewidth}P{0.48\linewidth}@{}}
\toprule
Artifact & Purpose and validation status\\\midrule
\endhead
\path{article/article.tex}, \path{article/article.pdf} & This article and
its compiled, visually checked PDF.\\
\path{code/independent_checks.py} & Executed independent mathematical
checks; requires SymPy.\\
\path{code/AuditSupport.wl} & Proposed nonintrusive Wolfram helpers;
not executed natively.\\
\path{patches/make_export_patch.py} & Source-pinned export containment
patch generator; Python transformation tests passed.\\
\path{tests/test_python_tools.py} & Seven executed Python-tool tests.\\
\path{tests/Baseline.wlt}, \path{tests/Support.wlt} & Unexecuted native
baseline/helper tests.\\
\path{tests/ProposedContracts.wlt} & Unexecuted desired contracts; not a
claimed passing upstream suite.\\
\path{tests/RunAudit.wls} & Native runner with explicit suite selection,
source file hash, and per-test result export.\\
\path{benchmarks/CompareBuiltins.wls} & Unexecuted native comparison
harness; no automated unmatched speedup claims.\\
\path{evidence/} & Executed independent results, native-execution status,
source inventory, coverage, findings, and provenance.\\
\path{README.md}, \path{UPSTREAM-LICENSE.txt} & Reproduction instructions,
limitations, and the license for included upstream excerpts.\\
\bottomrule
\end{longtable}

\clearpage
\begin{thebibliography}{99}
\addcontentsline{toc}{section}{References}
\bibitem{readme} Vladimir Reshetnikov and Asymptotic contributors.
Current package overview. \path{README.md}.
Pinned revision \code{07a9781212be}; accessed 9 September 2026.
\url{https://github.com/VladimirReshetnikov/Asymptotic/blob/07a9781212beb2eeb9ff16aa625b50ac27974078/README.md}.

\bibitem{core} Vladimir Reshetnikov and Asymptotic contributors.
Core implementation. \path{AsymptoticInverse/Kernel/AsymptoticInverse.wl}.
Pinned revision \code{07a9781212be}; accessed 9 September 2026.
\url{https://github.com/VladimirReshetnikov/Asymptotic/blob/07a9781212beb2eeb9ff16aa625b50ac27974078/AsymptoticInverse/Kernel/AsymptoticInverse.wl}.

\bibitem{seriesops} Vladimir Reshetnikov and Asymptotic contributors.
Series operations and replay refinement. \path{AsymptoticInverse/Kernel/SeriesOperations.wl}.
Pinned revision \code{07a9781212be}; accessed 9 September 2026.
\url{https://github.com/VladimirReshetnikov/Asymptotic/blob/07a9781212beb2eeb9ff16aa625b50ac27974078/AsymptoticInverse/Kernel/SeriesOperations.wl}.

\bibitem{seriesarith} Vladimir Reshetnikov and Asymptotic contributors.
Automatic arithmetic and held normalization. \path{AsymptoticInverse/Kernel/SeriesArithmetic.wl}.
Pinned revision \code{07a9781212be}; accessed 9 September 2026.
\url{https://github.com/VladimirReshetnikov/Asymptotic/blob/07a9781212beb2eeb9ff16aa625b50ac27974078/AsymptoticInverse/Kernel/SeriesArithmetic.wl}.

\bibitem{refinement} Vladimir Reshetnikov and Asymptotic contributors.
Retained inverse computation state. \path{AsymptoticInverse/Kernel/RefinementState.wl}.
Pinned revision \code{07a9781212be}; accessed 9 September 2026.
\url{https://github.com/VladimirReshetnikov/Asymptotic/blob/07a9781212beb2eeb9ff16aa625b50ac27974078/AsymptoticInverse/Kernel/RefinementState.wl}.

\bibitem{fourier} Vladimir Reshetnikov and Asymptotic contributors.
Fourier coefficient calculus. \path{AsymptoticInverse/Kernel/FourierCoefficients.wl}.
Pinned revision \code{07a9781212be}; accessed 9 September 2026.
\url{https://github.com/VladimirReshetnikov/Asymptotic/blob/07a9781212beb2eeb9ff16aa625b50ac27974078/AsymptoticInverse/Kernel/FourierCoefficients.wl}.

\bibitem{native} Vladimir Reshetnikov and Asymptotic contributors.
Native special-function series boundary. \path{AsymptoticInverse/Kernel/NativeSpecialFunctions.wl}.
Pinned revision \code{07a9781212be}; accessed 9 September 2026.
\url{https://github.com/VladimirReshetnikov/Asymptotic/blob/07a9781212beb2eeb9ff16aa625b50ac27974078/AsymptoticInverse/Kernel/NativeSpecialFunctions.wl}.

\bibitem{sourcecoords} Vladimir Reshetnikov and Asymptotic contributors.
Real source-coordinate transformations. \path{AsymptoticInverse/Kernel/SourceCoordinates.wl}.
Pinned revision \code{07a9781212be}; accessed 9 September 2026.
\url{https://github.com/VladimirReshetnikov/Asymptotic/blob/07a9781212beb2eeb9ff16aa625b50ac27974078/AsymptoticInverse/Kernel/SourceCoordinates.wl}.

\bibitem{certificates} Vladimir Reshetnikov and Asymptotic contributors.
Exact elementary inverse certificates. \path{AsymptoticInverse/Kernel/InverseCertificates.wl}.
Pinned revision \code{07a9781212be}; accessed 9 September 2026.
\url{https://github.com/VladimirReshetnikov/Asymptotic/blob/07a9781212beb2eeb9ff16aa625b50ac27974078/AsymptoticInverse/Kernel/InverseCertificates.wl}.

\bibitem{special} Vladimir Reshetnikov and Asymptotic contributors.
Explicit special-function adapters. \path{AsymptoticInverse/Kernel/SpecialFunctionAdapters.wl}.
Pinned revision \code{07a9781212be}; accessed 9 September 2026.
\url{https://github.com/VladimirReshetnikov/Asymptotic/blob/07a9781212beb2eeb9ff16aa625b50ac27974078/AsymptoticInverse/Kernel/SpecialFunctionAdapters.wl}.

\bibitem{gamma} Vladimir Reshetnikov and Asymptotic contributors.
Ordered Gamma/Barnes inverse construction. \path{AsymptoticInverse/Kernel/GammaInverse.wl}.
Pinned revision \code{07a9781212be}; accessed 9 September 2026.
\url{https://github.com/VladimirReshetnikov/Asymptotic/blob/07a9781212beb2eeb9ff16aa625b50ac27974078/AsymptoticInverse/Kernel/GammaInverse.wl}.

\bibitem{paclet} Vladimir Reshetnikov and Asymptotic contributors.
Paclet metadata. \path{AsymptoticInverse/PacletInfo.wl}.
Pinned revision \code{07a9781212be}; accessed 9 September 2026.
\url{https://github.com/VladimirReshetnikov/Asymptotic/blob/07a9781212beb2eeb9ff16aa625b50ac27974078/AsymptoticInverse/PacletInfo.wl}.

\bibitem{validation} Vladimir Reshetnikov and Asymptotic contributors.
Current and historical validation record. \path{validation/README.md}.
Pinned revision \code{07a9781212be}; accessed 9 September 2026.
\url{https://github.com/VladimirReshetnikov/Asymptotic/blob/07a9781212beb2eeb9ff16aa625b50ac27974078/validation/README.md}.

\bibitem{reviewmanifest} Vladimir Reshetnikov and Asymptotic contributors.
Historical review manifest. \path{validation/review-manifest.json}.
Pinned revision \code{07a9781212be}; accessed 9 September 2026.
\url{https://github.com/VladimirReshetnikov/Asymptotic/blob/07a9781212beb2eeb9ff16aa625b50ac27974078/validation/review-manifest.json}.

\bibitem{development} Vladimir Reshetnikov and Asymptotic contributors.
Development provenance and standalone packaging. \path{docs/development/README.md}.
Pinned revision \code{07a9781212be}; accessed 9 September 2026.
\url{https://github.com/VladimirReshetnikov/Asymptotic/blob/07a9781212beb2eeb9ff16aa625b50ac27974078/docs/development/README.md}.

\bibitem{ci} Vladimir Reshetnikov and Asymptotic contributors.
Standalone freshness CI workflow. \path{.github/workflows/standalone.yml}.
Pinned revision \code{07a9781212be}; accessed 9 September 2026.
\url{https://github.com/VladimirReshetnikov/Asymptotic/blob/07a9781212beb2eeb9ff16aa625b50ac27974078/.github/workflows/standalone.yml}.

\bibitem{runtests} Vladimir Reshetnikov and Asymptotic contributors.
Full-suite native runner. \path{AsymptoticInverse/Tests/RunTests.wl}.
Pinned revision \code{07a9781212be}; accessed 9 September 2026.
\url{https://github.com/VladimirReshetnikov/Asymptotic/blob/07a9781212beb2eeb9ff16aa625b50ac27974078/AsymptoticInverse/Tests/RunTests.wl}.

\bibitem{license} Vladimir Reshetnikov and Asymptotic contributors.
Repository license text. \path{LICENSE}.
Pinned revision \code{07a9781212be}; accessed 9 September 2026.
\url{https://github.com/VladimirReshetnikov/Asymptotic/blob/07a9781212beb2eeb9ff16aa625b50ac27974078/LICENSE}.

\bibitem{refinementrequests} Vladimir Reshetnikov and Asymptotic contributors.
Structured refinement requests. \path{AsymptoticInverse/Kernel/RefinementRequests.wl}.
Pinned revision \code{07a9781212be}; accessed 9 September 2026.
\url{https://github.com/VladimirReshetnikov/Asymptotic/blob/07a9781212beb2eeb9ff16aa625b50ac27974078/AsymptoticInverse/Kernel/RefinementRequests.wl}.

\bibitem{tree} Asymptotic repository tree at the reviewed revision.
\url{https://github.com/VladimirReshetnikov/Asymptotic/tree/07a9781212beb2eeb9ff16aa625b50ac27974078}.

\bibitem{asymptotic} Wolfram Research. \emph{Wolfram Language Documentation},
\code{Asymptotic}. Accessed 9 September 2026.
\url{https://reference.wolfram.com/language/ref/Asymptotic.html}.

\bibitem{asymptoticsolve} Wolfram Research. \emph{Wolfram Language Documentation},
\code{AsymptoticSolve}. Accessed 9 September 2026.
\url{https://reference.wolfram.com/language/ref/AsymptoticSolve.html}.

\bibitem{series} Wolfram Research. \emph{Wolfram Language Documentation},
\code{Series}. Accessed 9 September 2026.
\url{https://reference.wolfram.com/language/ref/Series.html}.

\bibitem{seriesdata} Wolfram Research. \emph{Wolfram Language Documentation},
\code{SeriesData}. Accessed 9 September 2026.
\url{https://reference.wolfram.com/language/ref/SeriesData.html}.

\bibitem{inverseseries} Wolfram Research. \emph{Wolfram Language Documentation},
\code{InverseSeries}. Accessed 9 September 2026.
\url{https://reference.wolfram.com/language/ref/InverseSeries.html}.

\bibitem{composeseries} Wolfram Research. \emph{Wolfram Language Documentation},
\code{ComposeSeries}. Accessed 9 September 2026.
\url{https://reference.wolfram.com/language/ref/ComposeSeries.html}.

\bibitem{seriescoefficient} Wolfram Research. \emph{Wolfram Language Documentation},
\code{SeriesCoefficient}. Accessed 9 September 2026.
\url{https://reference.wolfram.com/language/ref/SeriesCoefficient.html}.

\bibitem{productlog} Wolfram Research. \emph{Wolfram Language Documentation},
\code{ProductLog}. Accessed 9 September 2026.
\url{https://reference.wolfram.com/language/ref/ProductLog.html}.

\bibitem{asymptoticintegrate} Wolfram Research. \emph{Wolfram Language Documentation},
\code{AsymptoticIntegrate}. Accessed 9 September 2026.
\url{https://reference.wolfram.com/language/ref/AsymptoticIntegrate.html}.

\bibitem{discrete} Wolfram Research. \emph{Wolfram Language Documentation},
\code{DiscreteAsymptotic}. Accessed 9 September 2026.
\url{https://reference.wolfram.com/language/ref/DiscreteAsymptotic.html}.

\bibitem{asymptoticdsolve} Wolfram Research. \emph{Wolfram Language Documentation},
\code{AsymptoticDSolveValue}. Accessed 9 September 2026.
\url{https://reference.wolfram.com/language/ref/AsymptoticDSolveValue.html}.

\bibitem{asymptoticless} Wolfram Research. \emph{Wolfram Language Documentation},
\code{AsymptoticLess}. Accessed 9 September 2026.
\url{https://reference.wolfram.com/language/ref/AsymptoticLess.html}.

\bibitem{optionvalue} Wolfram Research. \emph{Wolfram Language Documentation},
\code{OptionValue}. Accessed 9 September 2026.
\url{https://reference.wolfram.com/language/ref/OptionValue.html}.

\end{thebibliography}
\end{document}
